% !TEX TS-program = pdflatex
\documentclass[11pt,a4paper,ragged2e,withhyper]{altacv}

\geometry{left=1.2cm,right=1.2cm,top=1.2cm,bottom=1.2cm,columnsep=1.2cm}

\usepackage{multicol}
\usepackage[utf8]{inputenc}
\usepackage[T1]{fontenc}

% If using pdflatex (no system font dependency):
\usepackage[rm]{roboto}
\usepackage[defaultsans]{lato}
\renewcommand{\familydefault}{\sfdefault}

\definecolor{RoyalBlue}{HTML}{2B4F81}
\definecolor{DarkGray}{HTML}{444444}
\colorlet{name}{RoyalBlue}
\colorlet{heading}{RoyalBlue}
\colorlet{accent}{RoyalBlue}
\colorlet{emphasis}{DarkGray}
\colorlet{body}{DarkGray}

\name{Eduardo Dueñez}
\tagline{Associate Professor of Mathematics at The University of Texas at San Antonio}
\personalinfo{
  \email{eduardo.duenez@utsa.edu}
  \location{San Antonio, TX, USA}
  \homepage{www.researchgate.net/profile/Eduardo-Duenez}
  \github{eduenez}
}

\begin{document}
\makecvheader

\cvsection{Education}
\cvevent{Ph.D. in Mathematics}{Princeton University}{2001}{Thesis advisor: Peter Sarnak.}
\cvevent{B.Sc. in Mathematics}{University of Guanajuato, Mexico}{1996}{}

\cvsection{Professional Employment}
\cvevent{Research Consultant}{Gibran.ai (AI startup; seed funding from Together Fund)}{2025 -- 2026}{Full-time research consulting, June-December 2025.}
\cvevent{Associate Professor}{The University of Texas at San Antonio}{2011 -- Present}{}
\cvevent{Assistant Professor}{The University of Texas at San Antonio}{2005 -- 2011}{}
\cvevent{Presidential Fellow}{The University of Texas at San Antonio}{2004 -- 2005}{}
\cvevent{J. J. Sylvester Assistant Professor}{The Johns Hopkins University}{2002 -- 2004}{}
\cvevent{Research Fellow}{American Institute of Mathematics, Palo Alto}{2001}{}

\cvsection{Research Awards \& Grants}
\cvachievement{\faTrophy}{NSF Grant: ``Model Theory and Ergodic Theorems'' (2015 -- 2019)}{Co-PI with José N. Iovino.}
\cvachievement{\faTrophy}{AMS/Tensor-SUMMA Grant (2007 -- 2008)}{Supported UTSA's Problem-Solving Seminar.}

\cvsection{Research Interests}
My research revolves around analysis in a broad sense. I have worked on random matrix theory (Dyson's “Threefold Way”; linear group ensembles) and number theory (zeros of L-functions); nowadays, my research is on the foundations of analysis and probability, as well as both theoretical and applied aspects of real-valued computing.

Recent and current research is on foundations of analysis from the perspective of real-valued and other continuous logics, including:
\begin{itemize}
  \item Measure, integration, probability, and information theory
  \item Real-valued computability (deterministic, randomized, and quantum)
  \item Statistical/machine learning and AI foundations
\end{itemize}

\cvsection{Refereed Publications}
\begin{enumerate}
\item E. Due{\~n}ez, A. S. Hamakiotes, S. J. Miller. \emph{Sums of Powers by {L}'{H}ospital's Rule.} Fibonacci Quart. \textbf{63}:3 (2025) 603--613.
\item Xavier Caicedo, Eduardo Due\~{n}ez, Jos\'e N. Iovino. \emph{Metastable convergence and logical compactness.} Beyond First Order Model Theory, Volume II, CRC Press, Boca Raton, FL (2023).
\item Eduardo Due\~{n}ez, Jos\'{e} N. Iovino. \emph{Model theory and metric convergence {II}: {A}verages of unitary polynomial actions.} Contemp. Math. \textbf{775} (2021) 85--114.
\item Eduardo Due\~{n}ez, Jos\'{e} N. Iovino. \emph{Model theory and metric convergence {I}: {M}etastability and dominated convergence.} Beyond First Order Model Theory, Volume I, CRC Press, Boca Raton, FL (2017) 131--187.
\item E. Due{\~n}ez, D. K. Huynh, J. P. Keating, S. J. Miller, N. C. Snaith. \emph{A Random Matrix Model for Elliptic Curve {$L$}-functions of Finite Conductor.} Journal of Physics A: Mathematical and Theoretical \textbf{45}:11 (2012) 115207 (32~pp).
\item E. Due{\~n}ez, D. K. Huynh, J. P. Keating, S. J. Miller, N. C. Snaith. \emph{The Lowest eigenvalue of {J}acobi random matrix ensembles and {P}ainlev{\'e} {VI}.} Journal of Physics A: Mathematical and Theoretical \textbf{43}:40 (2010) 405204 (27 pp).
\item E. Due{\~n}ez, D. W. Farmer, S. Froehlich, C. P. Hughes, F. Mezzadri, T. Phan. \emph{Roots of the Derivative of the {R}iemann Zeta Function and of Characteristic Polynomials.} Nonlinearity \textbf{23} (2010) 2599--2621.
\item E. Due{\~n}ez, S. J. Miller. \emph{The effect of convolving families of {$L$}-functions on the underlying group symmetries.} Proc. London Math. Soc. \textbf{99}:3 (2009) 787--820.
\item Stephanie Deacon, Eduardo Due{\~n}ez, Jos{\'e} N. Iovino. \emph{A public-key threshold cryptosystem based on residue rings.} J. Discrete Math. Sci. Cryptogr. \textbf{10}:4 (2007) 559--571.
\item E. Due{\~n}ez, S. J. Miller. \emph{The low lying zeros of a {$\rm GL(4)$} and a {$\rm GL(6)$} family of {$L$}-functions.} Compos. Math. \textbf{142}:6 (2006) 1403--1425.
\item E. Due{\~n}ez. \emph{``{H}arder'' ensembles of orthogonal matrices.} Experiment. Math. \textbf{15}:3 (2006) 273--277.
\item E. Due{\~n}ez, S. J. Miller, A. Roy, H. Straubing. \emph{Incomplete quadratic exponential sums in several variables.} J. Number Theory \textbf{116}:1 (2006) 168--199.
\item E. Due{\~n}ez. \emph{Random matrix ensembles associated to compact symmetric spaces.} Comm. Math. Phys. \textbf{244}:1 (2004) 29--61.
\end{enumerate}

\cvsection{Preprints and Work in Progress}
\begin{itemize}
\item E. Due{\~n}ez, J. Iovino, L. Salvetti Martinez, T. M. Wiederhold, Franklin Tall. \emph{Complexity of deep computations via topology of function spaces} (In preparation.) \href{https://doi.org/10.48550/arXiv.2601.00528}{arXiv:2601.00528}.
\item Samson Alva, E. Dueñez, J. N. Iovino, C. Walton. \emph{Deep equilibria: {E}xistence and computability} (To appear in Mathematical Structures in Computer Science.) \href{https://doi.org/10.48550/arXiv.2409.06064}{arXiv:2409.06064}.
\item Peter G. Casazza, Eduardo Due\~{n}ez, Jos\'e N. Iovino. \emph{Nondefinability of the spaces of Tsirelson and Schlumprecht} (Accepted.)
\item Eduardo Due\~{n}ez, Jos\'{e} N. Iovino. \emph{Analysis: The Study of Real-Valued and Random-Valued Structures} (Book in preparation.)
\end{itemize}

\cvsection{Academic Synergy and Mentorship}
\begin{itemize}
\item\emph{Since 2006.} Faculty mentor, UTSA Problem-Solving Club.
\item\emph{Since 2006.} UTSA Faculty Liaison to the Putnam Competition. (UTSA ranked top 8\% nationwide in 2015.)
\item\emph{2012.} Academic coach/mentor, AAAS Olympiad program. (Mentor/coach of 20 students from underrepresented communities in mathematics.)
\end{itemize}

\cvsection{Theses Supervised}
\subsection*{Master's Theses}
\begin{itemize}
  \item Zachery Sharon, UTSA (2011). \emph{An Explicit Proof of the Weak Finite Basis Theorem and Applications to Computing Ranks of Elliptic Curves}
  \item Elkin O. Quintero Vanegas, Universidad Nacional de Colombia (2011). \emph{La Fórmula del Número de Clases y Valores Especiales de Funciones-L: Desde Dirichlet y Dedekind hasta Stark}
\end{itemize}

\subsection*{Undergraduate Theses}
\begin{itemize}
  \item Julio C. Galindo López, UNAM (2013). \emph{Matrices Aleatorias y Funciones-L}
  \item Dal S. Yu, UTSA (2011). \emph{Representation Theory of Classical Compact Lie Groups}
\end{itemize}

\cvsection{Professional/Editorial Service and Distinctions}
\begin{itemize}
  \item Associate Editor, Mathematics Magazine (2016 -- 2019).
  \item Co-editor, Problems and Solutions column, Mathematics Magazine (2016 -- 2019).
  \item Co-organizer (with J. Iovino), session ``Topological Methods in Model Theory'', Meeting of the AMS Central Section, University of Texas at San Antonio (2024).
  \item Member, SACNAS Summer Leadership Institute (2015).
  \item Reviewer of manuscripts for Princeton University Press, Oxford University Press.
  \item Referee. Annals of Pure and Applied Logic, Bulletin of the London Mathematical Society, Communications in Mathematical Physics, Experimental Mathematics, Journal of Number Theory, Journal of Physics A, Mathematics Magazine, The Ramanujan Journal, Indian Journal of Mathematics, Random Matrices: Theory and Applications, and others.
  \item NSF panelist; proposal reviewer for CONACyT (Mexico).
  \item Judge of student abstracts and posters, SACNAS (2010 -- 2018).
\end{itemize}

\cvsection{Professional Memberships}
\begin{itemize}
  \item American Mathematical Society (AMS)
  \item Mathematical Association of America (MAA)
  \item Association for Symbolic Logic (ASL)
  \item Society for Advancement of Chicanos and Native Americans in Science (SACNAS)
\end{itemize}

\cvsection{Research Talks \& Short Courses}
{\footnotesize
\begin{itemize}
\item Special Session on Ergodic Theory and Topological Dynamics (ISQGD-SS14), International Society for Quantization, Topology and Dynamics, March 2026.
  \emph{Real-valued computations: a perspective via topology and model theory.}
\item Fields Institute (Toronto), Dec 2023.
  \emph{Integration of generalized functions in real-valued structures.}
\item CMS Winter Meeting (Toronto), Dec 2023.
  \emph{Structures of random variables and stability of Orlicz spaces.}
\item Seminario FAMAT/DEMAT, Universidad de Guanajuato, Feb 2021.
  \emph{Análisis y Lógica Continua.}
\item Universidad de Colima, Sep 2018.
  \emph{Convergencia metaestable y teoremas ergódicos desde la perspectiva de Lógica Continua.}
\item Casa Matemática Oaxaca (Mexico), June 2018.
  \emph{Convergencia metaestable y teoremas ergódicos desde la perspectiva de Lógica Continua.}
\item AMS Sectional Meeting, University of Central Florida (Orlando, FL), September 2017.
  \emph{Mean convergence of polynomial averages and continuous logic.}
\item Mathematical Congress of the Americas.  Montreal (Canada), July 2017.
  \emph{Ergodic theorems and metastability.}
\item Simposio Latinoamericano de Logica y Modelos. Puebla (Mexico), June 2017.
  \emph{Metastable convergence and ergodic theorems: a continuous logic viewpoint.}
\item Joint Mathematics Meeting (Atlanta), January 2017.
  \emph{Polynomial ergodic descent and uniformly metastable convergence.}
\item Joint Mathematics Meeting (Atlanta), January 2017.
  \emph{Uniformly metastable convergence in metric structures.}
\item National Security Agency, Ft. Meade, MD, October 2016.
  \emph{Metastable convergence of ergodic averages and continuous model theory.}
\item SACNAS National Conference. Long Beach, CA, October 2016.
  \emph{Random matrices and automorphic L-functions: Low-lying zeros and symmetry.}
\item MAA Southeastern Section Conference, University of Alabama at Birmingham, March 2016.
  \emph{The locus of median lines of Sierpinski's Triangle.}
\item Georgia Tech, February 2016.
  \emph{Convergence of ergodic averages and continuous model theory.}
\item Joint Mathematics Meeting. San Antonio, January 2015.
  \emph{The Mean Ergodic Theorem for Abelian Unitary Actions.}
\item SACNAS national conference. Los Angeles, October 2014.
  \emph{Families of polynomial mappings between groups.}
\item University of Texas at Austin, March 2014.
  \emph{Some arithmetic problems related to Sierpinski's Triangle.}
\item Mathematical Congress of the Americas. Guanajuato, Mexico, August 2013.
  \emph{A random matrix model for elliptic curve L-functions of finite conductor.}
\item Purdue University, March 2013.
  \emph{L-Functions and Random Matrices in Classical Compact Groups.}
\item SACNAS National Conference, October 2012.
  \emph{Random Matrices at the Intersection of Number Theory and Lie Groups.}
\item Encuentro conjunto RSME-SMM. Malaga, Spain. January 2012.
  \emph{Matrices aleatorias y ceros de funciones-L de twists cuadraticos de curvas elipticas.}
\item Universidad Pedagogica y Tecnologica de Tunja, Colombia, November 2011.
  \emph{Una invitacion a la teoria de representaciones.}
\item Universidad Nacional de Colombia, July 2011.
  \emph{Funciones Zeta y L: Aplicaciones clasicas y modernas en teoria de numeros.}
\item University of Colima, Mexico, November 2010.
  \emph{Funciones Zeta y L: Que son y para que sirven?}
\item U. T. Austin, October 2010.
  \emph{Zeros of L-functions and Random Matrices: Two Recent Applications.}
\item XLI Congreso de la Sociedad Matematica Mexicana, October 2008.
  \emph{La Zeta y las L's: Como se Metieron las Matrices Aleatorias a la Aritmetica?}
\item X CNTA, Waterloo, July 2008.
  \emph{Derivatives of Zeros of the Riemann Zeta-Function and of Characteristic Polynomials of Random matrices.}
\item Texas AM University, April 2008.
  \emph{Derivatives of Zeros of the Riemann Zeta-Function and of Characteristic Polynomials of Random matrices.}
\item 22nd Annual Workshop on Automorphic Forms and Related Topics, March 2008.
  \emph{Low-Lying Zeros of Modular (and Automorphic) L-functions.}
\item XL Congress of the Mexican Mathematical Society, UANL, October 2007.
  \emph{Condensacion de Zeros de Funciones-L de Curvas Elipticas.}
\item The Banff International Research Station (BIRS), July 2007.
  \emph{Finite-Conductor Models for Zeros of Elliptic Curves.}
\item Number Theory Fest at the University of Illinois at Urbana-Champaign, May 2007.
  \emph{Repulsion of Low Zeros in Families of Elliptic Curves.}
\item Texas AM University, December 2006.
  \emph{Convolution of Families of GL(2) L-functions.}
\item Instituto de Matemáticas, UNAM. Morelia (Mexico), August 2006.
  \emph{What Dyson Wanted to Tell Schrodinger: Random Matrices in Number Theory.}
\item IX CNTA, Vancouver, July 2006.
  \emph{Examples and Conjectures on Convolutions of Families of L-functions.}
\item Progress in Random Matrix Theory and Number Theory, University of Rochester, June 2006.
  \emph{Twists and Functorial Liftings of Families of L-functions.}
\item Joint Mathematics Meeting, San Antonio, January 2006.
  \emph{Symmetry in Twisted Families of L-functions.}
\item National Congress of the Mexican Mathematical Society, Mexico City, October 2005.
  \emph{Riemann, Dyson and Schrodinger: Random Matrices in Number Theory.}
\item The 2005 Conference on Automorphic Forms and Related Topics, University of North Texas, April 2005.
  \emph{Counterexample to a Conjecture of Keating and Snaith.}
\item The University of Texas at Austin, November 2004.
  \emph{Symmetry Beyond Root Numbers.}
\item The Pennsylvania State University, October 2004.
  \emph{Symmetric Space Ensembles with an Application to Elliptic Curves.}
\item University of Texas at San Antonio colloquium, September 2004.
  \emph{Random Matrices: From Quantum Physics to Number Theory.}
\item Isaac Newton Institute (U.K.), July 2004.
  \emph{Symmetry Beyond Root Numbers: a GL(6) example.}
\item Isaac Newton Institute (U.K.), June 2004.
  \emph{Symmetry Flipping in Families of L-functions.}
\item Centre de Recherches Mathematiques (University of Montreal, Canada), May 2004.
  \emph{Counterexample to a Conjecture of Keating and Snaith.}
\item American Mathematical Society Eastern Meeting, Lawrenceville, NJ, May 2004.
  \emph{The Symmetry Type of Two Families of L-functions.}
\item Augusta State University, Augusta, GA, February 2004.
  \emph{Statistical Properties of Energy Levels of Quantum Systems.}
\item The Ohio State University, October 2003.
  \emph{Random Matrices and Number Theory: A Happy Marriage?}
\item The West Coast Number Theory Conference, Asilomar, December 2001.
  \emph{Do the Cohen-Delaunay-Lenstra Heuristics apply to Quadratic Twists?}
\item Stanford University, November 2001.
  \emph{Random Matrix Ensembles associated to Compact Symmetric Spaces.}
\item Centre de Recherches Mathematiques (University of Montreal, Canada), August 2001.
  \emph{Random Matrix Ensembles associated to Compact Symmetric Spaces.}
\item The Johns Hopkins University, February 2001.
  \emph{Compact Symmetric Spaces and their associated Matrix Ensembles.}
\item The American Institute of Mathematics, May 2001.
  \emph{Breaking Universality in a Family of Circular Ensembles.}
\item The Institute for Advanced Study, February 2001.
  \emph{The Gaudin-Mehta Method and Classical Random Matrix Theory.}
\item Mathematisches Forschungsinstitut Oberwolfach, Germany, November 2000.
  \emph{Compact Symmetric Spaces and their associated Matrix Ensembles.}
\item National Congress of the Mexican Mathematical Society, Morelia, Mexico, November 1994.
  \emph{Div, Grad and Curl According to Lie.}
\end{itemize}
}

\cvsection{Survey and Outreach Talks \& Short Courses}
{\footnotesize
\begin{itemize}
\item\emph{From round-robin tournaments and networking events to combinatorial designs.} Morehouse College. Atlanta, GA, April 2018.
\item\emph{From round-robin tournaments and networking events to combinatorial designs.} Columbus State University, Columbus, GA, April 2016.
\item\emph{Introducción a la Teoría Clásica de Matrices Aleatorias.} 5-hour course at Centro de Investigación en Matemáticas, Merida, Mexico, April 2016.
\item\emph{Introducción a la Teoría Clásica de Matrices Aleatorias.} 5-hour course at Universidad Autónoma de Sinaloa, Mexico, October 2015.
\item\emph{El triángulo de Sierpinski, matrices modulares y aritmética.} SIDIM at Ponce, Puerto Rico. February 2014.
\item\emph{Signed Harmonic Series and Quadratic Number Fields.} Conference of the STMC (South Texas Mathematics Consortium). April 2013.
\item\emph{Matrices Aleatorias Gaussianas y Circulares.} 7-hour course at ``Escuela de Matrices Aleatorias'', CIMAT, Mexico, November 2012.
\item\emph{Minicurso de Matrices Aleatorias.} 7-hour course. ``Metodos Estocasticos en Sistemas Dinamicos''. CIMAT, Mexico, February 2009.
\item\emph{An Invitation to Diophantine Geometry.} XL Congress of the Mexican Mathematical Society, Monterrey, October 2007.
\item\emph{Random Matrices: A Meeting Point of Multiple Theories.} Trinity University, San Antonio, September 2005.
\item\emph{Three Lectures on Number Theory and Random Matrices.} The Ohio State University, August 2004.
\item\emph{Card Tricks, de Bruijn Sequences, and Hindu Music.} The McGraw Hill Center for Teaching and Learning (Princeton University), September 2000.
\item\emph{The Kinematic Method in Geometry.} UNAM, Mexico City, May 1994.
\end{itemize}
}

\cvsection{Teaching Experience}
Asterisks ``*'' indicate cross-listed (UG/GR) or graduate (GR) courses.
\begin{multicols}{2}
\begin{enumerate}[labelindent=0pt,labelwidth=2cm,leftmargin=!]
\item[Spring 2026] Mathematical Foundations of AI*
\item[Fall 2025] (Industrial consulting leave)
\item[Spring 2025] Foundations of Cryptography*
\item[Spring 2025] Modern Abstract Algebra
\item[Fall 2024] Foundations of Analysis
\item[Fall 2024] Modern Algebra
\item[Spring 2023] Foundations of Cryptography*
\item[Spring 2023] Abstract Algebra
\item[Fall 2022] Foundations of Mathematics
\item[Fall 2022] Algebra I*
\item[Spring 2022] Foundations of Cryptography*
\item[Spring 2022] Algebra and Number Systems
\item[Fall 2021] Linear Algebra
\item[Fall 2021] Algebra and Number Systems
\item[Spring 2021] Foundations of Cryptography*
\item[Spring 2021] Algebra and Number Systems
\item[Fall 2020] Foundations of Mathematics
\item[Fall 2020] Foundations of Analysis
\item[Spring 2020] Foundations of Cryptography*
\item[Fall 2019] Linear Algebra
\item[Fall 2019] Topics in Number Theory*
\item[Spring 2019] Foundations of Cryptography*
\item[Spring 2019] Linear Algebra
\item[Fall 2018] Linear Algebra
\item[Fall 2018] Calculus II
\item[Spring 2018] Abstract Algebra II
\item[Spring 2018] Foundations of Mathematics
\item[Fall 2017] Abstract Algebra I
\item[Fall 2017] Precalculus
\item[Spring 2017] Foundations of Cryptography*
\item[Spring 2017] Foundations of Mathematics
\item[Fall 2016] Foundations of Mathematics
\item[Fall 2016] Calculus I
\item[Spring 2016] (Research leave)
\item[Fall 2015] (Research leave)
\item[Spring 2015] Foundations of Cryptography*
\item[Spring 2015] Linear Algebra
\item[Fall 2014] Foundations of Mathematics
\item[Fall 2014] Calculus I
\item[Fall 2014] Problem-Solving Seminar
\item[Spring 2014] Foundations of Analysis
\item[Spring 2014] Differential Equations I
\item[Fall 2013] Calculus II
\item[Fall 2013] Problem-Solving Seminar
\item[Spring 2013] Foundations of Analysis
\item[Spring 2013] Complex Variables
\item[Fall 2012] Calculus I
\item[Fall 2012] Multilinear Algebra
\item[Spring 2012] Modern Abstract Algebra
\item[Spring 2012] Calculus I
\item[Fall 2011] Number Theory
\item[Fall 2011] Problem-Solving Seminar
\item[Spring 2011] Linear Algebra
\item[Spring 2011] Problem-Solving Seminar
\item[Fall 2010] Calculus II
\item[Fall 2010] Problem-Solving Seminar
\item[Spring 2010] Calculus II
\item[Spring 2010] Problem-Solving Seminar
\item[Fall 2009] Number Theory
\item[Fall 2009] Problem-Solving Seminar
\item[Spring 2009] Elliptic-Curve Cryptography*
\item[Spring 2009] Problem-Solving Seminar
\item[Fall 2008] Calculus III
\item[Fall 2008] Linear Algebra
\item[Fall 2008] Problem-Solving Seminar
\item[Spring 2008] Calculus II
\item[Spring 2008] Problem-Solving Seminar
\item[Fall 2007] Problem-Solving Seminar
\item[Fall 2007] Foundations of Mathematics
\item[Spring 2007] Problem-Solving Seminar
\item[Spring 2007] Foundations of Analysis
\item[Spring 2007] Linear Algebra
\item[Fall 2006] General Topology I*
\item[Fall 2006] Calculus III
\item[Spring 2006] Foundations of Analysis
\item[Fall 2005] Complex Variables I*
\item[Fall 2005] Freshman Seminar
\item[Summer 2005] Calculus II
\item[Spring 2005] Modern Abstract Algebra
\item[Fall 2004] Modern Algebra
\item[Spring 2004] Random Matrices with Applications
\item[Spring 2004] Fourier Analysis
\item[Fall 2003] Complex Analysis with Applications
\item[Spring 2003] Calculus II for Biology and Medicine
\item[Fall 2002] Calculus I for Biology and Medicine
\item[Spring 2002] Linear Algebra
\item[Fall 1999] Vector calculus
\end{enumerate}
\end{multicols}

\end{document}